\section{Worked compilation patterns}\label{sec:cases}

\subsection{Case study I: least witness and finite minimum}

Let \(g(i,x)\) be a primitive-recursive natural-valued function and let \(N>0\).  We want a public
index \(j<N\) minimizing \(g(j,x)\).

Define the bounded predicate
\[
 P(j,x,N)\;:\Longleftrightarrow\;j<N\land
   \forall i<N\,[g(j,x)\le g(i,x)].
\]
By compositional bounded reflection, \(P\) has a bit.  Existence can be proved without classical
choice by a trace that scans \(0,\ldots,N-1\) while maintaining the best index.  The state is
\((t,j)\), initially \((1,0)\), and the step replaces \(j\) by \(t\) when
\(g(t,x)<g(j,x)\).  The invariant says
\[
 j<t\quad\land\quad\forall i<t\;g(j,x)\le g(i,x).
\]
The bounded universal is reflected to a bit so the induction formula is quantifier free.  At
\(t=N\) the invariant gives \(P(j,x,N)\).

If the public theorem only needs the minimum value, the index can be sealed and the value shown
independent of the selected minimizer.  If the theorem needs a canonical index, use the least index
satisfying \(P\), supplied by bounded choice.

\paragraph{Schedule.}
A compact schedule has nodes:

\begin{enumerate}
\item defining equations for the scan trace;
\item reflection of the bounded invariant;
\item base invariant;
\item preservation;
\item quantifier-free induction;
\item final minimum theorem;
\item optional least-minimizer naming.
\end{enumerate}

There are no residuals: all data are primitive recursive and the finite nonempty domain is explicit.

\subsection{Case study II: a bounded game strategy}

Consider
\[
 \forall x<N\;\exists y<M(x)\;R(x,y,q),
\]
where \(R\) is primitive recursive and the intended public result is the existence of a response
function for the fixed finite domain.

For each \(x<N\), define
\[
 w(x)=\min\{y<M(x):R(x,y,q)\},
\]
using the arity-polymorphic bounded-choice theorem.  The premise required at arbitrary \(x\) is the
object formula \(x<N\to\exists y<M(x)\,R(x,y,q)\).  Universal elimination supplies it under the
hypothesis \(x<N\).  Bounded choice yields \(w(x)<M(x)\) and \(R(x,w(x),q)\).  Universal
introduction closes \(x\).

The function \(w\) is primitive recursive because the search bound and verifier are.  This is a
legitimate finite choice function; no unbounded choice principle is used.

If the public theorem merely asserts the original bounded sentence, constructing \(w\) is optional.
One can eliminate each existential locally.  The named strategy is justified only because the
bounds are explicit and the family is finite.

For a second alternation
\[
 \forall x<N\;\exists y<M(x)\;\forall z<K(x,y)\;\exists u<L(x,y,z)\;R(x,y,z,u),
\]
compile the inner existential first, reflect the inner bounded universal, and then search for an outer
\(y\) whose reflected bit is one.  Meta-level recursion on the fixed quantifier tree produces the
strategy terms and their proofs.

\subsection{Case study III: bisection for a monotone inverse}

Suppose \(F\) is represented on rationals by exact arithmetic and has a proved monotonicity theorem.
Given rational \(a<b\), target \(y\), and signs
\[
 F(a)\le y\le F(b),
\]
a fixed number \(p\) of bisection steps produces an interval \([a_p,b_p]\) of width
\((b-a)/2^p\) that still brackets the target.

Encode dyadic rationals by integer numerator and exponent.  The state contains endpoints and a
failure bit.  At each step compute the midpoint \(m\), compare \(F(m)\) with \(y\), and keep the
appropriate half.  The quantifier-free invariant is:

\begin{enumerate}
\item the endpoints have the expected dyadic denominator;
\item \(a_t\le b_t\);
\item \(F(a_t)\le y\le F(b_t)\);
\item the width identity \(2^t(b_t-a_t)=b-a\);
\item the failure bit is zero.
\end{enumerate}

The base is immediate.  Preservation uses exact midpoint arithmetic and monotonicity.  One
quantifier-free induction proves the invariant at \(p\).  The public output can be the interval or
its midpoint with an explicit error bound.

An exact inverse value is not required.  If an external theorem gives a separation gap ensuring that
only one algebraic branch intersects the interval, a later schedule node recovers the branch label.
Without that residual, the interval is the maximal constructive result.

\subsection{Case study IV: height-bounded Diophantine enumeration}\label{subsec:height}

Let \(S(q,z)\) be a primitive-recursive solution predicate for codes \(z\), and let
\(h(z)\) be a primitive-recursive height.  Assume a domain-specific producer theorem
\[
 B(q,H):\quad \forall z\,[S(q,z)\to h(z)\le H].
\]
Also assume a complete coding bound \(C(q,H)\) with an object proof
\[
 h(z)\le H\to z<C(q,H).
\]
The second condition is often elementary once the code convention is fixed; the first may be the
deep theorem.

Define the list bit for every \(z<C(q,H)\) by filtering with \(S(q,z)\).  A list code can be produced
by a primitive-recursive scan, or the theorem can use the characteristic bit directly.

\begin{theorem}[Conditional complete-list theorem]\label{thm:complete-list}
Under \(B(q,H)\), iPRA proves:
\begin{enumerate}
\item every listed code is a solution;
\item every solution code occurs in the bounded list;
\item equality and any finite observable can be decided by bounded search over the list.
\end{enumerate}
\end{theorem}

\begin{proof}
Soundness is built into the filter: membership implies the solution bit is one, hence \(S\) by
reflection.  For completeness, let \(z\) satisfy \(S(q,z)\).  The producer theorem gives
\(h(z)\le H\), and the coding theorem gives \(z<C(q,H)\).  The scan visits \(z\), its solution bit is
one, and the list invariant proves membership.  All list operations are bounded and therefore
primitive recursive.
\end{proof}

This theorem resolves the logical packaging of Faltings/Siegel-style packets: once a genuine height
producer is supplied, complete enumeration and all subsequent finite reasoning compile uniformly.
It does not manufacture the height theorem.  The schedule endpoint is
\[
 B(q,H),\Gamma\vdash_{\PRAAx}\mathsf{CompleteList}(q,H,L).
\]
A packet that lacks \(B\) should publish exactly this conditional theorem rather than claiming an
unconditional complete list.

\subsection{Case study V: a certificate cloud and a gap gate}

Let \(C_p(q,c)\) be a validity predicate for precision-\(p\) certificates, and let \(a(c)\) be a
rational observation.  Assume:
\[
 \exists c\,C_p(q,c),
 \qquad
 C_p(q,c)\land C_p(q,d)\to |a(c)-a(d)|\le2^{-p}.
\]
The first witness may be entirely local.  Select it existentially, output the rational interval
\[
 I_p(c)=[a(c)-2^{-p},a(c)+2^{-p}],
\]
and seal \(c\).  Any other valid certificate has observation in this interval, so the interval is a
certificate-independent public cloud up to the stated enlargement.

Now suppose labels \(\ell<K\) have reference values \(v_\ell\), and a residual gap theorem states
\[
 \ell\ne m\to |v_\ell-v_m|>4\cdot2^{-p},
\]
while every valid certificate is within \(2^{-p}\) of the reference value of some label.  Then the
label is unique.  A bounded search through \(K\) labels finds it, and the uniqueness proof shows that
the least found label is the intended one.

The architecture separates two theorems:

\begin{enumerate}
\item unconditionally, produce a small cloud or stable observation;
\item conditionally on the gap, produce an exact label.
\end{enumerate}

This prevents approximation accuracy from being collapsed into constructive exactness.

\subsection{Case study VI: a finite cover argument}

Suppose a theorem about a compact-looking domain uses only a finite set of test points once a mesh
\(\delta\) is fixed.  Let \(N(q,\delta)\) be a supplied net and let \(\mathsf{Covers}\) be a
primitive-recursive verifier that every coded basic region is within \(\delta\) of a net point.  A
uniform continuity modulus transports a property from the net to the whole finite-use query.

The schedule has residual nodes for the net-completeness theorem and the modulus if they are not
already checked.  The remaining nodes are bounded:

\begin{enumerate}
\item verify the property at every net point;
\item reflect the bounded universal;
\item use the modulus and cover certificate to transport to an arbitrary query point;
\item seal the chosen nearby net point.
\end{enumerate}

No unrestricted compactness principle occurs.  The public theorem states exactly the finite cover and
modulus assumptions consumed.

\subsection{Case study VII: recurrence identities}

For sequences such as Fibonacci, Lucas, Catalan, or Bell numbers, the repository already contains
many checked recurrence and bounded-sum derivations.  \TAPIR{} adds value when a new identity combines
several libraries.

For example, to prove a fixed polynomial identity in \(F_n,F_{n+1},L_n\):

\begin{enumerate}
\item import the recurrence equations as closed plans;
\item choose a quantifier-free polynomial invariant in the paired state
      \((F_n,F_{n+1},L_n,L_{n+1})\);
\item prove the base by normalization;
\item prove preservation by rewriting with recurrence equations and the checked ring tool;
\item compile one induction node and specialize at the requested \(n\).
\end{enumerate}

If an identity involves a bounded convolution, the sum reflection module supplies the recursion and
its equations.  The schedule compiler combines these local derivations without introducing a new
logic or a new semantic model.

\subsection{What these examples show}

The examples share one architecture:

\begin{enumerate}
\item finite code or fixed precision;
\item primitive-recursive checker or trace;
\item bounded reflection and quantifier-free induction;
\item local sealing of noncanonical data;
\item named witnesses only under explicit bounds;
\item residual hypotheses for the genuinely deep producer theorems;
\item schedule compilation to a literal object sequent.
\end{enumerate}

The framework is broad because these proof moves recur, not because iPRA is being asked to contain an
unrestricted universe of mathematics.
